\documentclass[11pt]{article}

\usepackage[letterpaper,margin=1in]{geometry}
\usepackage[T1]{fontenc}
\usepackage[utf8]{inputenc}
\usepackage{lmodern}
\usepackage{microtype}
\usepackage{amsmath,amssymb,mathtools,bm}
\usepackage{xcolor}
\usepackage{enumitem}
\usepackage{hyperref}
\usepackage[most]{tcolorbox}
\usepackage{fancyhdr}
\usepackage{parskip}

\definecolor{qftblue}{HTML}{245A8D}
\definecolor{qftgreen}{HTML}{2F6B4F}
\definecolor{qftred}{HTML}{A23B3B}
\definecolor{qftgray}{HTML}{F4F6F8}

\hypersetup{
  colorlinks=true,
  linkcolor=qftblue,
  urlcolor=qftblue,
  pdftitle={PHYS 662 Homework Set 1},
  pdfauthor={Nima Lashkari}
}

\pagestyle{fancy}
\fancyhf{}
\lhead{PHYS 662: Quantum Field Theory I}
\rhead{Homework Set 1}
\cfoot{\thepage}
\setlength{\headheight}{14pt}

\setlist[enumerate]{leftmargin=2.2em,itemsep=0.55em,topsep=0.5em}
\newlist{parts}{enumerate}{2}
\setlist[parts,1]{label=\textbf{(\alph*)},leftmargin=2.3em,itemsep=0.7em}
\setlist[parts,2]{label=\textbf{(\roman*)},leftmargin=2.4em,itemsep=0.45em}

\newcommand{\dd}{\mathrm{d}}
\newcommand{\ii}{\mathrm{i}}
\newcommand{\R}{\mathbb{R}}
\newcommand{\cH}{\mathcal{H}}

\newtcolorbox{conventionbox}[1][]{
  enhanced,
  breakable,
  colback=qftgray,
  colframe=qftblue,
  boxrule=0.8pt,
  arc=1.5mm,
  title={#1},
  fonttitle=\bfseries,
  coltitle=white
}

\newcommand{\problem}[2]{%
  \clearpage
  \section*{\color{qftblue}Problem #1: #2}
  \addcontentsline{toc}{section}{Problem #1: #2}
}

\begin{document}

\begin{center}
  {\scshape\Large PHYS 662 --- Quantum Field Theory I\par}
  \vspace{0.25em}
  {\Huge\bfseries\color{qftblue} Homework Set 1\par}
  \vspace{0.25em}
  {\large Lectures 1--4 \quad Fall 2026\par}
\end{center}

\vspace{0.5em}
\noindent
\textbf{Name:} \rule{0.44\textwidth}{0.4pt}
\hfill
\textbf{Due:} \rule{0.24\textwidth}{0.4pt}

\begin{conventionbox}[Instructions and conventions]
Show the reasoning leading to each result and state any additional
assumptions. Unless a problem says otherwise, use periodic boundary
conditions and natural units \(\hbar=c=k_B=1\). Fourier-transform
normalizations must be stated and used consistently.
\end{conventionbox}

\noindent
The four problems are weighted equally.

\tableofcontents

\problem{1}{Translation-invariant quadratic lattices}

Consider \(N\) identical degrees of freedom \((q_j,p_j)\) on a periodic
one-dimensional lattice, with \(j=0,\ldots,N-1\) understood modulo \(N\).
Let
\begin{equation}
  H=
  \sum_{j=0}^{N-1}\frac{p_j^2}{2m}
  +\frac12\sum_{j,\ell=0}^{N-1}q_j K_{j\ell}q_\ell,
  \qquad
  K_{j\ell}=K(j-\ell).
  \label{eq:generalLattice}
\end{equation}
Assume that \(K_{j\ell}\) is real, symmetric, and positive semidefinite.

\begin{parts}
  \item Define a discrete translation by \(n\) sites and show explicitly that
  Eq.~\eqref{eq:generalLattice} is invariant under it. Distinguish this
  discrete translation of lattice sites from a uniform displacement
  \(q_j\mapsto q_j+c\).

  \item Introduce the discrete Fourier transform
  \begin{equation}
    q_j=\frac{1}{\sqrt N}\sum_k Q_k\,e^{\ii kaj},
    \qquad
    p_j=\frac{1}{\sqrt N}\sum_k P_k\,e^{\ii kaj},
    \qquad
    k=\frac{2\pi n}{Na},\quad n=0,\ldots,N-1.
    \label{eq:DFT}
  \end{equation}
  State the reality conditions on \(Q_k\) and \(P_k\). Determine the normal
  modes and show that their frequencies satisfy
  \begin{equation}
    \omega^2(k)=\frac{\widetilde K(k)}{m},
    \qquad
    \widetilde K(k)=\sum_r K(r)e^{-\ii kar}.
  \end{equation}
  Explain why the assumptions on \(K\) make \(\widetilde K(k)\) real and
  nonnegative.

  \item What additional condition on \(K\) makes the Hamiltonian invariant
  under \(q_j\mapsto q_j+c\)? Relate it to the \(k=0\) normal mode.

  \item Now consider a square lattice of spacing \(a\), with one real scalar
  displacement \(q_{\bm r}\) at each site and identical nearest-neighbor
  springs:
  \begin{equation}
    H=\sum_{\bm r}\frac{p_{\bm r}^{2}}{2m}
    +\frac{\kappa}{2}
    \sum_{\bm r}\sum_{\alpha=1}^{2}
    \left(q_{\bm r+a\hat e_\alpha}-q_{\bm r}\right)^2,
    \qquad \kappa>0.
    \label{eq:squareLattice}
  \end{equation}
  Find the plane-wave normal modes and derive
  \begin{equation}
    \omega^2(\bm k)
    =\frac{4\kappa}{m}
    \left[
      \sin^2\left(\frac{k_xa}{2}\right)
      +\sin^2\left(\frac{k_ya}{2}\right)
    \right].
    \label{eq:squareDispersion}
  \end{equation}
  Obtain the long-wavelength dispersion and identify the propagation speed.
\end{parts}

\problem{2}{The one-dimensional ferromagnetic Heisenberg model}

Consider a periodic chain of quantum spins of length \(s\),
\begin{equation}
  H=-J\sum_{j=0}^{N-1}\bm S_j\cdot\bm S_{j+1},
  \qquad J>0,
  \label{eq:Heisenberg}
\end{equation}
with \(\hbar=1\) and
\begin{equation}
  [S_i^\alpha,S_j^\beta]
  =\ii\delta_{ij}\epsilon^{\alpha\beta\gamma}S_i^\gamma,
  \qquad
  \bm S_i^{2}=s(s+1).
  \label{eq:spinAlgebra}
\end{equation}

\begin{parts}
  \item Use the Heisenberg equation \(\dot{\bm S}_i=\ii[H,\bm S_i]\) to show
  that
  \begin{equation}
    \dot{\bm S}_i
    =J\bm S_i\times
    \left(\bm S_{i-1}+\bm S_{i+1}\right).
    \label{eq:spinEOM}
  \end{equation}
  Check directly that the total spin
  \(\bm S_{\rm tot}=\sum_i\bm S_i\) is conserved.

  \item Take the large-\(s\) classical limit and write
  \(\bm S_i=s\bm n(x_i,t)\), where \(x_i=ia\) and
  \(|\bm n|=1\). Expand Eq.~\eqref{eq:spinEOM} through order \(a^2\) and show
  that the long-wavelength field obeys
  \begin{equation}
    \partial_t\bm n
    =Js a^2\,\bm n\times\partial_x^2\bm n
    +\mathcal O(a^4\partial_x^4).
    \label{eq:LandauLifshitz}
  \end{equation}
  Verify that this evolution preserves \(|\bm n|=1\).

  \item Linearize about the ferromagnetic ground state
  \(\bm n=\hat{\bm z}\). With
  \(\psi=n_x+\ii n_y\), show that
  \begin{equation}
    \ii\partial_t\psi=-Js a^2\partial_x^2\psi.
  \end{equation}
  Determine the long-wavelength magnon dispersion and compare it with the
  exact lattice result
  \(\omega(k)=2Js(1-\cos ka)\).
\end{parts}

\problem{3}{Gaussian \texorpdfstring{\(O(N)\)}{O(N)} model and correlation length}

Let \(\bm\phi=(\phi^1,\ldots,\phi^N)\) be a classical field in \(d\)
Euclidean dimensions with free-energy functional
\begin{equation}
  \mathcal F[\bm\phi]
  =\frac12\int \dd^d x\,
  \left[
    \partial_i\phi^a\partial_i\phi^a
    +r\,\phi^a\phi^a
  \right],
  \qquad r>0,
  \label{eq:GaussianON}
\end{equation}
and thermal weight \(Z^{-1}e^{-\mathcal F/T}\). Repeated indices are summed.

\begin{parts}
  \item Show that Eq.~\eqref{eq:GaussianON} is invariant under spatial
  rotations and under the independent internal transformation
  \(\phi^a(x)\mapsto U^a{}_b\phi^b(x)\), \(U\in O(N)\). Under which group is
  \(\phi^a\) a scalar, and under which group is it a vector?

  \item Fourier transform the field and diagonalize \(\mathcal F\). Use the
  resulting Gaussian integral to derive
  \begin{equation}
    \langle\phi^a(\bm k)\phi^b(\bm k')\rangle
    =(2\pi)^d\delta^{ab}\delta^{(d)}(\bm k+\bm k')
    \frac{T}{\bm k^2+r}.
    \label{eq:GaussianCorrelator}
  \end{equation}

  \item Show that the correlation length is
  \(\xi=r^{-1/2}\). In \(d=3\), evaluate the real-space two-point function and
  show that
  \begin{equation}
    \langle\phi^a(\bm x)\phi^b(\bm 0)\rangle
    =\delta^{ab}\frac{T}{4\pi|\bm x|}
    e^{-|\bm x|/\xi}.
  \end{equation}

  \item Define the susceptibility per component by
  \begin{equation}
    \chi=\frac1N\int\dd^d x\,
    \langle\phi^a(\bm x)\phi^a(\bm 0)\rangle.
  \end{equation}
  Compute \(\chi\), and explain what happens to \(\xi\) and \(\chi\) as
  \(r\to0^+\). Why does this limit make the separation between microscopic
  and long-distance scales especially large?
\end{parts}

\problem{4}{Fock space and spatial versus internal symmetry}

Consider \(N\) real quantum scalar fields in \(d\) spatial dimensions,
\begin{equation}
  H=\frac12\int\dd^d x\,
  \left[
    \pi^a\pi^a
    +\partial_i\phi^a\partial_i\phi^a
    +m^2\phi^a\phi^a
  \right],
  \label{eq:multiscalarH}
\end{equation}
with equal-time commutators
\begin{equation}
  [\phi^a(t,\bm x),\pi^b(t,\bm y)]
  =\ii\delta^{ab}\delta^{(d)}(\bm x-\bm y),
  \qquad
  [\phi^a,\phi^b]=[\pi^a,\pi^b]=0.
  \label{eq:canonicalFieldCCR}
\end{equation}

\begin{parts}
  \item Show that the Hamiltonian is invariant under
  \begin{align}
    \text{spatial rotations: }&
    \phi_R^a(t,\bm x)=\phi^a(t,R^{-1}\bm x),
    &&R\in SO(d),\notag\\
    \text{internal rotations: }&
    \phi_U^a(t,\bm x)=U^a{}_b\phi^b(t,\bm x),
    &&U\in O(N),
  \end{align}
  with analogous transformations for \(\pi^a\). Show that the two actions
  commute. Explain precisely why \(\phi^a\) is a spatial scalar but an
  internal vector.

  \item Let \(\omega_{\bm k}=\sqrt{\bm k^2+m^2}\) and use
  \begin{equation}
    \phi^a(t,\bm x)=
    \int\frac{\dd^d k}{(2\pi)^d}
    \frac{1}{\sqrt{2\omega_{\bm k}}}
    \left[
      a^a_{\bm k}e^{-\ii\omega_{\bm k}t+\ii\bm k\cdot\bm x}
      +a^{a\dagger}_{\bm k}
       e^{\ii\omega_{\bm k}t-\ii\bm k\cdot\bm x}
    \right].
    \label{eq:multiscalarExpansion}
  \end{equation}
  Determine the corresponding expansion of \(\pi^a=\dot\phi^a\) and derive
  \begin{equation}
    [a^a_{\bm k},a^{b\dagger}_{\bm k'}]
    =(2\pi)^d\delta^{ab}\delta^{(d)}(\bm k-\bm k').
    \label{eq:modeCCR}
  \end{equation}
  Verify the remaining equal-time commutators.

  \item Show that a one-particle state
  \(a^{a\dagger}_{\bm k}|0\rangle\) transforms as a vector of \(O(N)\), while
  its momentum is rotated by \(SO(d)\). State the transformation law for
  \(a^{a\dagger}_{\bm k}\) under each group.

  \item Construct, up to normalization, an \(O(N)\)-singlet two-particle state
  with total momentum zero. Explain why there is no nonzero one-particle
  singlet in the defining vector representation for \(N>1\).

  \item Use Eq.~\eqref{eq:multiscalarExpansion} to verify that
  \begin{equation}
    [\phi^a(t,\bm x),\phi^b(t,\bm y)]=0.
  \end{equation}
  Explain how this equal-time statement fits into the broader locality
  requirement discussed in Lecture 2.
\end{parts}

\end{document}
